
encore int\`egre sur \(S\), \`a fibres de dimension \(2\)\,; en vertu de 4.3, on peut remplacer \(X\) par \(Y\). Enfin, en vertu de 4.3 et 3.10 (utilisant la r\'esolution des singularit\'es des surfaces), quitte \`a remplacer \(S\) par un rev\^etement plat d'un ouvert non-vide de \(S\), on peut remplacer \(X\) par un \(S\)-sch\'ema lisse.

Maintenant, gr\^ace \`a 4.6 (i) \(\Longleftrightarrow\) (iii) et au th\'eor\`eme de finitude et de sp\'ecialisation des groupes de cohomologie \'etale (SGA 4 XVI 2.2), 5.1 r\'esulte de la proposition suivante.

\textbf{Proposition 5.2.} \emph{Soient \(k\) un corps alg\'ebriquement clos et \(X\) un \(k\)-sch\'ema propre et lisse. Alors le groupe \(C^p_{\rm num}(X)\) des cycles alg\'ebriques de codimension \(p\) sur \(X\), modulo l'\'equivalence num\'erique, est libre de rang \(\le b^{2p}\), le nombre de Betti.}

En effet, si l'on oublie l'action des racines de 1, la cohomologie \'etale \(H^*(X,\Q_\ell)=_{\rm dfn}\left[\varprojlim H^*(X,\Z/\ell^\nu\Z)\right]\otimes_{\Z_\ell}\Q_\ell\) o\`u \((\ell,\operatorname{car}(k))=1\), devient une cohomologie de Weil~: elle satisfait \`a la dualit\'e de Poincar\'e et \`a la formule de K\"Unneth\,; et elle re\c{c}oit un homomorphisme fonctoriel, non nul de l'anneau des cycles alg\'ebriques modulo l'\'equivalence rationnelle \(C^*_{\rm rat}(X)\) (voir X 7.13).

Soient \(a_1,\ldots,a_m\in C^{r-p}_{\rm rat}(X)\), o\`u \(r=\dim(X)\), des \'el\'ements qui s'en vont sur une base du \(\Q_\ell\)-espace vectoriel dans \(H^{2(r-p)}(X,\Q_\ell)\), engendr\'e par les cycles alg\'ebriques. Consid\'erons l'homomorphisme
\[
   \alpha:C^p_{\rm rat}(X)\longrightarrow \Z^m
\]
d\'efini par \(\alpha(b)=(\langle b.a_1\rangle,\ldots,\langle b,a_m\rangle)\), o\`u \(\langle\ \rangle\) d\'esigne le nombre d'intersection. Alors on voit facilement que l'image de \(\alpha\) est \(C^p_{\rm num}(X)\).

